\documentclass{report}
\usepackage{amsmath,amssymb}
\setlength{\parindent}{0mm}
\setlength{\parskip}{1em}
\begin{document}
\begin{center}
\rule{6in}{1pt} \
{\large Miguel Aguilo \\
{\bf Solving Large-Scale Inverse Problems in Elasticity Using Sequential Quadratic Programming}}

Optimization and Uncertainty Quantification \\ P O Box 5800 \\ MS 1318 \\ Sandia National Laboratories \\ Albuquerque \\ NM \\ USA
\\
{\tt maguilo@sandia.gov}\\
Denis Ridzal\\
Joseph Young\end{center}

In this work we use a matrix-free trust-region sequential
quadratic programming (SQP) algorithm to solve inverse
problems governed by the partial differential equations
(PDEs) of linear and nonlinear elasticity.
SQP methods are based on the successive minimization of a
quadratic model of the Lagrangian functional subject to a
linear Taylor model of the PDE constraints.
We derive first-order and second-order derivative operators
that are necessary for an accurate and efficient
implementation of these models in the context of elasticity
equations.

The linear systems arising in our approach are related to
convex quadratic programs that involve the model of the PDE
constraints, but not the model of the Lagrangian functional.
Such linear systems can be solved efficiently using
preconditioners based on available iterative PDE solvers.
The effectiveness of our approach is demonstrated through
numerical examples in linear elastodynamics and nonlinear
elastostatics.


\end{document}
